\begin{document}

\begin{frame}
\maketitle
\end{frame}

\begin{frame}[fragile]{Goals}
\begin{itemize}
    \item Minimality of the core language
        \begin{itemize}
            \item Only 3 special forms: \linline|quote|, and \linline|macroexpand| and \linline|cur-env|
            \item Implement everything else in \lang itself
        \end{itemize}
    \item Ease of implementing partial evaluation
    \item Expressiveness
        \begin{itemize}
            \item Full LISP-like macros: implement flow control, DSLs, custom syntax
            \item Control over continuations: implement error handling, state machines, async execution
        \end{itemize}
\end{itemize}
\end{frame}

\begin{frame}[fragile]{\lang: Features}
\begin{itemize}
    \item Lisp-like
        \begin{itemize}
            \item Closures
            \item Macros
        \end{itemize}
    \item Precise control over compile-time vs runtime execution
        \begin{itemize}
            \item Macro evaluation has ``runtime'' semantics: the macro body can make use of values defined in the runtime environment
            \item Constant data, including macro execution, may be evaluated at compile time during PE
        \end{itemize}
\end{itemize}
\end{frame}

\begin{frame}[fragile]{Example code in \lang}
\begin{minted}{lisp}
(module
  (doc "this module calculates factorials")
  (exports main)
  (imports
    (builtin
      (mul add <= macroexpand
        lambda quote))
    ("core/prelude.l" (if const fn)))
  (defs
    (!const main (!lambda () (fact 5)))

    (!fn (fact x)
      (!if (<= x 0)
        1
        (mul x (fact (add x -1)))))))
\end{minted}
\end{frame}

\begin{frame}[fragile]{Builtin special forms in \lang}
\begin{itemize}
  \item \linline|cur-env| -- introspect the environment at the callsite
  \item \linline|quote| -- return the given expression without evaluating
  \item \linline|macroexpand| -- apply given macro to quoted arguments and then evaluate result
  \begin{itemize}
    \item Syntax sugar is available (next slide)
  \end{itemize}
\end{itemize}
\end{frame}

\begin{frame}[fragile]{Requirements for bootstrapping}
Functions for manipulating basic data are assumed to be included in the initial environment:
\begin{itemize}
\item Lists: \linline|cons|, \linline|car|, \linline|cdr|, \linline|null?|
\item Tests: \linline|symbol?|, \linline|pair?|, \linline|func?|, \linline|sym-eq?|
\item Booleans: \tinline|#t|, \tinline|#f|, \linline|bool-to-k|
\item Functional: \linline|lambda|, \linline|eval|, \linline|apply|
\item Extrinsic: \linline|read-source|
\end{itemize}
\end{frame}

\begin{frame}[fragile]{Macro execution in \lang}
\begin{itemize}
    \item This code:
      \begin{minted}{lisp}
        (!if (= x 5)
          "x is 5"
          "x is not 5")
      \end{minted}
    \item Is syntax sugar for this code:
      \begin{minted}{lisp}
        (eval (cur-env) (if (quote (= x 5))
          (quote "x is 5")
          (quote "x is not 5")))
      \end{minted}
    \item Macros are just functions that operate on code.
\end{itemize}
\end{frame}

\begin{frame}[fragile]{Defining the building blocks of a module system}
Bootstrapping can be done in a few tiers:
\begin{itemize}
    \item \linline|list|, \linline|if|
    \item The simple \linline|Y| combinator (if not available)
    \item \linline|map|
    \item \linline|let|
    \item The polyvariate \linline|Y| combinator (if not available)
    \item \linline|letrec|
    \item Operations with key-value lists
    \item Module parsing
    \item Imports
\end{itemize}
\end{frame}

\begin{frame}[fragile]{Bootstrapping specific constructs}
\begin{minted}{lisp}
  ; definition of helper "expand-lambda"
  (!lambda (args body)
    (cons macroexpand
      (cons lambda (cons args (cons body ())))))

  ; definition of "map" without use of letrec
  (Y (!lambda (map)
      (!lambda (f l)
        (!if (null? l)
          l
          (!if (pair? l)
            (cons (f (car l)) (map f (cdr l)))
            (make-fail (list 'not-a-list l)))))))))

  ; definition of "let"
  (!lambda (letlist body)
    (cons
      (expand-lambda (map car letlist) body)
      (map cadr letlist)))))
\end{minted}
\end{frame}
\begin{frame}[fragile]{Core concepts}
\begin{itemize}
    \item Modules are converted to \linline|letrec| definitions that allow mutual recursion
    \item The result of evaluating these definitions is a key-value list of exported items
    \item Importing modules works by first evaluating them and then making their exports available in the new environment
\end{itemize}
\end{frame}

\begin{frame}[fragile]{Future research}
\begin{itemize}
    \item Partial evaluation of the module import process
      \begin{itemize}
        \item If \linline|read-source| can be evaluated as a constant during PE,
        a module import declaration can be specialised to the key-value list
        resulting from its evaluation
      \end{itemize}
    \item Additional features
      \begin{itemize}
        \item Declaration of constants in addition to functions
        \item Module structure verification
      \end{itemize}
\end{itemize}
\end{frame}

\begin{frame}[fragile,plain]
\begin{center}
\Huge Thank you!
\end{center}
\end{frame}
\end{document}
